This survey examines reinforcement learning broadly construed across economics, bridging two intellectual traditions that study sequential decision-making. The economics tradition begins with a formal model of agents, institutions, preferences and equilibrium, and works from there to observed choices and data; the reinforcement learning tradition begins with a specified environment and computes optimal behavior. Both traditions solve variants of the Bellman equation, and their convergence is the organizing theme of this survey. The curse of dimensionality limits exact dynamic programming, and a growing class of economic problems resists the simplification necessary for dynamic programming to work. Reinforcement learning algorithms offer a natural, sample-based extension that trades convergence guarantees for computational scalability. The survey proceeds through algorithms, theory, deep RL empirics, structural estimation, macroeconomic models, strategic games, bandits, offline reinforcement learning, reinforcement learning from human feedback, causal inference for and from reinforcement learning, distributional, robust, and constrained methods, world models, and evidence-gated field deployments. The emphasis throughout is on applying reinforcement learning to economic models and applied decision problems, not on formal treatment alone, and the computational chapters include working simulations demonstrating the method's leverage relative to classical alternatives. The survey also examines the practical vulnerabilities of these algorithms, noting their brittleness, sample inefficiency, sensitivity to hyperparameters, and the absence of global convergence guarantees outside of tabular settings. The successes of reinforcement learning remain strictly bounded by these constraints. The target audience is economists curious about reinforcement learning and reinforcement learning researchers curious about economics. A companion survey \citep{RustRawat2026} covers the inverse problem of inferring preferences from observed behavior. All simulation code is publicly available.\footnote{\url{https://github.com/rawatpranjal/survey-of-reinforcement-learning-in-economics}}
